\documentclass[11pt]{article}
\usepackage{amsmath}
\usepackage{graphicx}
\newcommand{\method}{\textsc{Freddie}}
\newcommand*{\vect}[1]{\mathbf{#1}}
\def\eg{e.g.,\xspace}
% A full-line comment that must disappear.
\title{A Study of Structured Translation}
\author{Ada Lovelace \and Alan Turing}
\date{}
\begin{document}
\maketitle
\begin{abstract}
We present \method, a system that translates papers while keeping their structure intact. It reaches 95\% accuracy on 10-12$\times$ longer inputs.
\end{abstract}
\begin{comment}
This block is a comment environment and must be removed.
\end{comment}
\section{Introduction}\label{sec:intro}
Machine translation of scientific papers must preserve equations such as $E=mc^2$ and references.
Prior work \cite{smith2020} translates text only; we instead translate \LaTeX{} sources directly.

\section{Method}\label{sec:method}
Our approach splits the source into units.   % trailing comment
Each unit is translated with its neighbors as context, see Section~\ref{sec:intro} and \cite{smith2020,doe2021}.
\begin{figure}[t]
  \centering
  \includegraphics[width=0.8\linewidth]{figures/overview.pdf}
  \caption{Overview of the pipeline. Units flow from parsing to reconstruction.}
  \label{fig:overview}
\end{figure}
\begin{itemize}
  \item First, we \emph{parse} the project.
  \item Second, we translate every unit\footnote{Concurrently, with bounded parallelism.}.
\end{itemize}
\begin{equation}
  \vect{y} = W \vect{x} + b
\end{equation}
\subsection{Validation}
We validate commands, placeholders and delimiters; failing units are retried up to three times.
\begin{table}[h]
\caption{Results on the benchmark.}
\begin{tabular}{lc}
Model & Score \\
Ours & 0.95 \\
\end{tabular}
\end{table}
\section{Conclusion}
The structure survives translation, as shown in Figure~\ref{fig:overview}.
\section*{Appendix}
Additional experiments are reported here. We ran each configuration five times and report the mean.

\bibliographystyle{plain}
\bibliography{refs}
\end{document}
